\section{Notation}
\label{sec:notation}

Throughout this work, plain italic type is used for scalars and scalar functions, for example the wavenumber $k$. Bold italic type is reserved for elements of $\R^d$ or $\C^d$ with $d \in \{2, 3\}$ and for functions with values in these sets, such as the electric field intensity $\V{E}$. Bold upright type denotes finite arrays of coefficients, such as vectors and matrices, as opposed to the vectors of physical space; an example is the impedance matrix $\M{Z}$. Calligraphic type is used for operators on function spaces, including functionals; for example, the Calder\'on identity~\cite{calderon1963boundary} reads $\OP{T}^2 = -\OP{I} / 4 + \OP{K}^2$. Number sets, already used above, are typeset in blackboard bold.
